\section{Attack Catalogue}
\label{sec:catalogue}

Table~\ref{tab:catalogue} lists the tactics we study. Each row is a transformation with one interface, one intensity parameter, and one implementation in the repository. The systems we execute are prompted rewrites, a local masked-LM filler, and deterministic text rules; several rows are operationalisations of a published attack rather than the published system. The row-by-row difference is in the released catalogue document: discourse and document rewrite prompt a model rather than running DIPPER, masked infilling masks a seeded sample rather than the highest-self-information tokens, synonym substitution takes candidates from a prompted model rather than Three Bricks' substitution model, and the pivot rows use one model for every leg. The search is non-adaptive: no tactic is trained against the detector. The families differ in what they perturb: paraphrase and regeneration replace the token stream through a model, back-translation replaces it through a translation system, and token edits touch individual positions. The control row is included so that a reader can see the difference between removing an edit-based carrier and attacking a statistical mark.

\begin{table*}[t]
\centering
\footnotesize
\caption{Attack catalogue. Cost is the expected loss of semantic fidelity at the intensity used in the main experiments: low, medium, or high. Source names the closest published relative, not necessarily the system executed; \S\ref{sec:catalogue} states the operational difference per row. Each row maps to a module under \texttt{efface/attacks/}.}
\label{tab:catalogue}
\begin{tabular}{@{}>{\raggedright\arraybackslash}p{3.4cm} >{\raggedright\arraybackslash}p{2.6cm} >{\raggedright\arraybackslash}p{3.5cm} l >{\raggedright\arraybackslash}p{3.3cm}@{}}
\toprule
Tactic & Family & What it changes & Cost & Closest relative \\
\midrule
Instruction paraphrase      & Paraphrase      & Sentence wording, kept close to source & Low    & \citet{kgwreliability} \\
Discourse paraphrase        & Paraphrase      & Paragraph structure and wording & Medium & \citet{dipper} \\
Masked-LM infilling         & Paraphrase      & Individual spans, local context only & Low & \citet{sira} \\
Style normalisation         & Paraphrase      & Contractions, casing, punctuation, filler & Low & \citet{markmywords} \\
Document rewrite            & Paraphrase      & Whole document, fresh context per chunk & High & \citet{dipper} \\
Single-pivot back-translation & Back-translation & Token stream via one pivot language & Medium & \citet{markmywords} \\
Multi-hop back-translation  & Back-translation & Token stream via two or more pivots & High & \citet{cwra} \\
Cross-lingual generation    & Back-translation & Response generated in the pivot, then translated & High & \citet{cwra} \\
Synonym substitution        & Token edits     & A fraction of content words & Low & \citet{threebricks} \\
Random word edits           & Token edits     & Swap, insert, or delete at random positions & Low & \citet{markmywords} \\
Typo and noise injection    & Token edits     & Character-level noise and misspellings & Low & \citet{markmywords} \\
Non-origin regeneration     & Regeneration    & The entire token stream, from a different model & High & \citet{cwra} \\
\midrule
Unicode scrub (control)     & Control         & Invisible and zero-width characters only & None & \citet{markmywords} \\
\bottomrule
\end{tabular}
\end{table*}

\paragraph{Paraphrase.}
Paraphrase replaces the writer's word choices while keeping the propositional content. The family spans a range of intensity. A style pass changes surface form and leaves most content words in place. Instruction-prompted paraphrase rewrites sentence by sentence and typically keeps most entities and numbers. DIPPER-style discourse paraphrase conditions on the surrounding passage and may reorder or merge sentences \citep{dipper}. Masked-LM infilling replaces a fraction of spans with predictions from a masked language model, which keeps the rest of the token stream intact \citep{sira}. Document rewriting discards local context and can drift furthest from the source, which is why it sits at the expensive end.

\paragraph{Back-translation.}
A pivot round trip translates the document into a pivot language and back. The output is fluent in the source language and usually preserves propositions, while the token sequence is the translator's rather than the generator's. \citet{cwra} show that watermark strength is not preserved across languages and that the round trip removes a large share of the signal. Multi-hop chains repeat the round trip through several pivots. Cross-lingual generation differs in what happens in the middle of the chain: the source is translated into the pivot, the attacker model rewrites the translated document in the pivot language, and the rewrite is translated back. That variant moves the watermark closer to its source of failure, because the middle step replaces the token stream in a language the watermark was never applied to.

\paragraph{Token edits.}
Token edits touch a bounded fraction of positions. Synonym substitution preserves meaning to the extent that the substitution model is reliable \citep{threebricks}. Random swaps, insertions, and deletions are weaker per edit but cheap to apply at scale \citep{markmywords}. Typo and noise injection perturbs characters, which changes tokenisation and therefore the context used to key the watermark. These tactics are the cheapest per document. A later full-stream replacement overwrites them, so the search treats them as last steps rather than as prefixes a rewrite then discards (\S\ref{sec:search}).

\paragraph{Regeneration.}
Regeneration asks a different model to produce a new document from the prompt, with the original available or withheld. It replaces every token and therefore removes the mark most completely, at the highest cost to fidelity. We include it as the upper bound on removal and the lower bound on preservation.

\paragraph{Control.}
The Unicode scrub removes zero-width characters, bidirectional controls, and exotic spaces. These carriers belong to edit-based watermarking, where the mark is inserted into the text as characters. The control row quantifies how much of an observed removal is attributable to that channel, and it should show no effect on a sampling-time mark.
